This appendix provides supplementary construction details and diagnostics for the methodology developed in Chapter~\ref{chap:methodology}. It documents the sample-construction procedures, feature definitions and diagnostics, preprocessing and selection rules, chronological sample allocation, and predictive-model specifications. The material supports the reproducibility of the empirical design while keeping the main chapter focused on its essential methodological choices.

\section{Sample Construction}
\label{app:sample_construction}

Filing selection is performed separately for each manager and reporting quarter. The SEC distinguishes holdings reports from notice reports through the EDGAR submission type and the cover-page report type \parencite[FAQs 20 and 33]{sec2023form13f}. For this analysis, a filing is eligible only if it was submitted by the common information cutoff and both classifications identify it as a holdings report. Notice filings, filings with conflicting classifications, and filings submitted after the cutoff are excluded. Among the remaining original filings and restatements, the base filing is selected by ordering candidates by amendment number, filing date, and accession number and retaining the final observation. Earlier bases are treated as superseded. A later amendment is retained as an addition only when its amendment number is greater than that of the selected base. This rule produces one reproducible set of filings for each eligible manager-quarter while using only information available by the prediction date.

For security mapping, nine-character CUSIPs are normalized and matched exactly to CRSP CUSIP histories whose validity intervals contain the relevant reporting date. A position is classified as matched only when this procedure identifies one unique PERMNO; observations with no candidate or multiple candidates are classified as unmatched or ambiguous, respectively. The holdings panel subsequently retains matched securities classified by CRSP as common equity and aggregates their reported positions first at the manager-security level and then into stock-quarter ownership measures.

Table~\ref{tab:sample_construction} summarizes the construction of the final modeling sample. Observations that lack either usable historical market features or a usable forward target remain in the full panel with explicit eligibility flags rather than being deleted. Overall, approximately 184,000 observations, or 89\% of the full panel, enter the three modeling samples, indicating that the eligibility and purging rules preserve broad coverage while enforcing the point-in-time design.

\begin{table}[H]
    \centering
    \caption{Stock-quarter panel and modeling-sample construction}
    \label{tab:sample_construction}
    \begin{tabular}{lr}
        \toprule
        Characteristic & Value \\
        \midrule
        Coverage & 2013 Q2--2026 Q1 \\
        Reporting quarters & 52 \\
        Distinct securities & 7,232 \\
        Full stock-quarter observations & 206,868 \\
        Usable predictors and target & 192,109 (92.87\%) \\
        Final modeling observations & 184,018 (88.95\%) \\
        \bottomrule
    \end{tabular}
\end{table}


\section{Downside-Risk Target Statistics}
\label{app:target_statistics}

Table~\ref{tab:target_statistics} summarizes the distributions of the primary 63-day maximum-drawdown target and the four alternative downside-risk definitions used in the robustness analysis. Statistics are calculated over usable observations from the common modeling universe, applying the corresponding usability rule to each target. All distributional statistics except skewness are reported as percentages.

\begin{table}[!htbp]
    \centering
    \scriptsize
    \setlength{\tabcolsep}{4pt}
    \caption{Distributional statistics for downside-risk targets}
    \label{tab:target_statistics}
    \begin{tabular}{lrrrrrrrr}
        \toprule
        Target definition & Mean & SD & P1 & P10 & Median & P90 & P99 & Skewness \\
        \midrule
        63-day maximum drawdown & 21.52\% & 16.22\% & 0.98\% & 6.24\% & 16.73\% & 44.20\% & 77.41\% & 1.52 \\
        63-day downside volatility & 33.62\% & 27.61\% & 2.58\% & 12.34\% & 25.92\% & 63.08\% & 132.38\% & 5.74 \\
        63-day worst five-day loss & 13.65\% & 10.97\% & 0.60\% & 4.39\% & 10.54\% & 26.48\% & 57.83\% & 2.38 \\
        21-day maximum drawdown & 12.03\% & 11.15\% & 0.55\% & 2.89\% & 8.50\% & 25.66\% & 55.71\% & 2.33 \\
        126-day maximum drawdown & 29.64\% & 20.06\% & 1.42\% & 9.31\% & 24.01\% & 59.44\% & 89.08\% & 1.07 \\
        \bottomrule
    \end{tabular}
\end{table}


\section{Market Feature Definitions and Diagnostics}
\label{app:market_features}

Table~\ref{tab:market-feature-definitions} gives the exact construction of the nine market predictors. Let $t$ denote the final trading day strictly before the information cutoff, $D_k$ the $k$ trading days ending on $t$, $r_{i,d}$ the daily return of stock $i$, and $r_{M,d}$ the CRSP total-market return. The complete feature window contains the final 252 trading days before the cutoff. Only regular-way observations with an active trading-status flag enter the calculations.

\begingroup
\scriptsize
\setlength{\tabcolsep}{4pt}
\renewcommand{\arraystretch}{1.15}
\begin{longtable}{>{\raggedright\arraybackslash}p{0.22\textwidth}>{\raggedright\arraybackslash}p{0.48\textwidth}>{\raggedright\arraybackslash}p{0.20\textwidth}}
        \caption{Exact definitions of the market predictors}
        \label{tab:market-feature-definitions}                                                                                                                                                                         \\
        \toprule
        Feature & Exact construction                                                                                                                                                         & Window and requirements \\
        \midrule
        \endfirsthead
        \toprule
        Feature & Exact construction                                                                                                                                                         & Window and requirements \\
        \midrule
        \endhead
        \bottomrule
        \endfoot
        \nolinkurl{market_cap_usd}
                & Most recent eligible CRSP daily market capitalization observed before the cutoff; the CRSP value is multiplied by 1,000 to express it in U.S. dollars.
                & Latest observation in the 252-day window.                                                                                                                                                            \\
        \nolinkurl{average_daily_dollar_volume_63d}
                & $|D_{63}|^{-1}\sum_{d\in D_{63}}|P_{i,d}|V_{i,d}$, where $P_{i,d}$ and $V_{i,d}$ are the CRSP daily price and volume.
                & 63 trading days.                                                                                                                                                                                     \\
        \nolinkurl{return_21d}
                & $\prod_{d\in D_{21}}(1+r_{i,d})-1$.
                & 21 trading days.                                                                                                                                                                                     \\
        \nolinkurl{momentum_252_21d}
                & Cumulative return from the beginning of the 252-day feature window through the day immediately preceding the final 21 trading days: $\prod_{d=t-251}^{t-21}(1+r_{i,d})-1$.
                & 231 trading days, with the most recent 21 days omitted.                                                                                                                                              \\
        \nolinkurl{volatility_63d}
                & Sample standard deviation of $r_{i,d}$ multiplied by $\sqrt{252}$.
                & 63 trading days; annualized.                                                                                                                                                                         \\
        \nolinkurl{downside_volatility_63d}
                & $\sqrt{252\,|D_{63}|^{-1}\sum_{d\in D_{63}}\min(r_{i,d},0)^2}$.
                & 63 trading days; annualized around a zero-return threshold.                                                                                                                                          \\
        \nolinkurl{market_beta_252d}
                & Ordinary-least-squares slope from regressing $r_{i,d}$ on the contemporaneous CRSP total-market return $r_{M,d}$.
                & 252 trading days; at least 202 paired observations.                                                                                                                                                  \\
        \nolinkurl{maximum_drawdown_252d}
                & Positive magnitude of the largest peak-to-trough decline in cumulative stock wealth, calculated by the same convention as the forward target.
                & 252 trading days.                                                                                                                                                                                    \\
        \nolinkurl{negative_return_share_63d}
                & $|D_{63}|^{-1}\sum_{d\in D_{63}}\mathbf{1}\{r_{i,d}<0\}$.
                & 63 trading days.                                                                                                                                                                                     \\
\end{longtable}
\endgroup

\noindent
A stock-quarter's market features are considered usable only when the 63-day window contains at least 50 observed returns, the 252-day window contains at least 202 observed returns, beta has at least 202 paired stock-market observations, and the latest price and market capitalization are positive. Figure~\ref{fig:market-feature-distributions} summarizes the training-sample distributions of the nine market predictors, with values clipped at their 1st and 99th percentiles for display only.

\begin{figure}[H]
        \centering
        \begin{subfigure}[t]{0.31\textwidth}
                \centering
                \includegraphics[width=\linewidth]{04_00_market_feature_distribution_01_market_cap_usd.pdf}
        \end{subfigure}
        \hfill
        \begin{subfigure}[t]{0.31\textwidth}
                \centering
                \includegraphics[width=\linewidth]{04_00_market_feature_distribution_02_average_daily_dollar_volume_63d.pdf}
        \end{subfigure}
        \hfill
        \begin{subfigure}[t]{0.31\textwidth}
                \centering
                \includegraphics[width=\linewidth]{04_00_market_feature_distribution_03_return_21d.pdf}
        \end{subfigure}
        \par\medskip
        \begin{subfigure}[t]{0.31\textwidth}
                \centering
                \includegraphics[width=\linewidth]{04_00_market_feature_distribution_04_momentum_252_21d.pdf}
        \end{subfigure}
        \hfill
        \begin{subfigure}[t]{0.31\textwidth}
                \centering
                \includegraphics[width=\linewidth]{04_00_market_feature_distribution_05_volatility_63d.pdf}
        \end{subfigure}
        \hfill
        \begin{subfigure}[t]{0.31\textwidth}
                \centering
                \includegraphics[width=\linewidth]{04_00_market_feature_distribution_06_downside_volatility_63d.pdf}
        \end{subfigure}
        \par\medskip
        \begin{subfigure}[t]{0.31\textwidth}
                \centering
                \includegraphics[width=\linewidth]{04_00_market_feature_distribution_07_market_beta_252d.pdf}
        \end{subfigure}
        \hfill
        \begin{subfigure}[t]{0.31\textwidth}
                \centering
                \includegraphics[width=\linewidth]{04_00_market_feature_distribution_08_maximum_drawdown_252d.pdf}
        \end{subfigure}
        \hfill
        \begin{subfigure}[t]{0.31\textwidth}
                \centering
                \includegraphics[width=\linewidth]{04_00_market_feature_distribution_09_negative_return_share_63d.pdf}
        \end{subfigure}
        \caption[Training-sample distributions of market predictors]{Training-sample distributions of the nine market predictors. Values are clipped at their 1st and 99th percentiles for display only.}
        \label{fig:market-feature-distributions}
\end{figure}

\noindent
Figure~\ref{fig:market-temporal-stability} tracks the quarterly medians of the market predictors over the training period. Their evolution reinforces the use of chronological validation and fitting-sample preprocessing rather than parameters estimated once from the full panel.

\begin{figure}[H]
        \centering
        \begin{subfigure}[t]{0.31\textwidth}
                \centering
                \includegraphics[width=\linewidth]{04_00_market_temporal_stability_01_market_cap_usd.pdf}
        \end{subfigure}
        \hfill
        \begin{subfigure}[t]{0.31\textwidth}
                \centering
                \includegraphics[width=\linewidth]{04_00_market_temporal_stability_02_average_daily_dollar_volume_63d.pdf}
        \end{subfigure}
        \hfill
        \begin{subfigure}[t]{0.31\textwidth}
                \centering
                \includegraphics[width=\linewidth]{04_00_market_temporal_stability_03_return_21d.pdf}
        \end{subfigure}
        \par\medskip
        \begin{subfigure}[t]{0.31\textwidth}
                \centering
                \includegraphics[width=\linewidth]{04_00_market_temporal_stability_04_momentum_252_21d.pdf}
        \end{subfigure}
        \hfill
        \begin{subfigure}[t]{0.31\textwidth}
                \centering
                \includegraphics[width=\linewidth]{04_00_market_temporal_stability_05_volatility_63d.pdf}
        \end{subfigure}
        \hfill
        \begin{subfigure}[t]{0.31\textwidth}
                \centering
                \includegraphics[width=\linewidth]{04_00_market_temporal_stability_06_downside_volatility_63d.pdf}
        \end{subfigure}
        \par\medskip
        \begin{subfigure}[t]{0.31\textwidth}
                \centering
                \includegraphics[width=\linewidth]{04_00_market_temporal_stability_07_market_beta_252d.pdf}
        \end{subfigure}
        \hfill
        \begin{subfigure}[t]{0.31\textwidth}
                \centering
                \includegraphics[width=\linewidth]{04_00_market_temporal_stability_08_maximum_drawdown_252d.pdf}
        \end{subfigure}
        \hfill
        \begin{subfigure}[t]{0.31\textwidth}
                \centering
                \includegraphics[width=\linewidth]{04_00_market_temporal_stability_09_negative_return_share_63d.pdf}
        \end{subfigure}
        \caption[Temporal stability of market predictors]{Quarterly medians of the market predictors across the training sample.}
        \label{fig:market-temporal-stability}
\end{figure}

\section{Ownership Feature Definitions and Diagnostics}
\label{app:crowding_measures}

Let $v_{m,i,q}$ denote manager $m$'s reported position value in stock $i$ during quarter $q$, $V_{i,q}=\sum_m v_{m,i,q}$ the stock's total reported institutional value, $P_{m,q}=\sum_j v_{m,j,q}$ manager $m$'s total reported portfolio value, $p_{m,i,q}=v_{m,i,q}/P_{m,q}$ the position's portfolio weight, and $M_q$ the number of reporting managers represented in the quarterly network. Table~\ref{tab:ownership-feature-definitions} defines the eleven conventional ownership predictors.

\begingroup
\scriptsize
\setlength{\tabcolsep}{4pt}
\renewcommand{\arraystretch}{1.15}
\begin{longtable}{>{\raggedright\arraybackslash}p{0.23\textwidth}>{\raggedright\arraybackslash}p{0.51\textwidth}>{\raggedright\arraybackslash}p{0.16\textwidth}}
        \caption{Exact definitions of the conventional ownership predictors}
        \label{tab:ownership-feature-definitions}                                                                                                                                              \\
        \toprule
        Feature & Exact construction                                                                                                                                       & Construction rule \\
        \midrule
        \endfirsthead
        \toprule
        Feature & Exact construction                                                                                                                                       & Construction rule \\
        \midrule
        \endhead
        \bottomrule
        \endfoot
        \nolinkurl{institutional_holder_count}
                & Number of distinct reporting managers with a retained position in the stock.
                & Calculated in levels.                                                                                                                                                        \\
        \nolinkurl{total_reported_value_usd}
                & $V_{i,q}=\sum_m v_{m,i,q}$.
                & Calculated in levels.                                                                                                                                                        \\
        \nolinkurl{reported_ownership_hhi}
                & $\sum_m (v_{m,i,q}/V_{i,q})^2$; concentration is measured within the stock's reported institutional ownership, not relative to total shares outstanding.
                & Used in levels.                                                                                                                                                              \\
        \nolinkurl{top_five_reported_ownership_share}
                & Sum of $v_{m,i,q}/V_{i,q}$ across the five managers with the largest reported position values in the stock.
                & Used in levels.                                                                                                                                                              \\
        \nolinkurl{mean_owner_portfolio_weight}
                & Unweighted mean of $p_{m,i,q}$ across the stock's reporting owners.
                & Used in levels.                                                                                                                                                              \\
        \nolinkurl{maximum_owner_portfolio_weight}
                & Maximum value of $p_{m,i,q}$ across the stock's reporting owners.
                & Used in levels.                                                                                                                                                              \\
        \nolinkurl{manager_holder_share}
                & Institutional holder count divided by $M_q$, the total number of reporting managers represented in quarter $q$.
                & Used in levels.                                                                                                                                                              \\
        \nolinkurl{institutional_holder_count_change_percent}
                & $100(H_{i,q}-H_{i,q-1})/H_{i,q-1}$, where $H_{i,q}$ is the institutional holder count.
                & Defined only for consecutive observed quarters.                                                                                                                              \\
        \nolinkurl{total_reported_value_change_percent}
                & $100(V_{i,q}-V_{i,q-1})/V_{i,q-1}$.
                & Defined only for consecutive observed quarters.                                                                                                                              \\
        \nolinkurl{reported_ownership_hhi_change}
                & $\mathrm{HHI}_{i,q}-\mathrm{HHI}_{i,q-1}$, an absolute rather than percentage change.
                & Defined only for consecutive observed quarters.                                                                                                                              \\
        \nolinkurl{top_five_reported_ownership_share_change}
                & Difference between the current and preceding quarter's top-five reported-ownership share.
                & Defined only for consecutive observed quarters.                                                                                                                              \\
\end{longtable}
\endgroup

\noindent
The level features require at least one retained manager position and have no rolling minimum-observation requirement. Change features require observations in consecutive reporting quarters. Missing change measures remain distinguishable at construction and are addressed only within the model preprocessing pipeline.

Figure~\ref{fig:ownership-feature-distributions} summarizes the training-sample distributions of the eleven conventional ownership predictors. Display-only clipping does not alter the values used in the analysis.

\begin{figure}[H]
        \centering
        \begin{subfigure}[t]{0.31\textwidth}
                \centering
                \includegraphics[width=\linewidth]{04_00_ownership_feature_distribution_01_institutional_holder_count.pdf}
        \end{subfigure}
        \hfill
        \begin{subfigure}[t]{0.31\textwidth}
                \centering
                \includegraphics[width=\linewidth]{04_00_ownership_feature_distribution_02_total_reported_value_usd.pdf}
        \end{subfigure}
        \hfill
        \begin{subfigure}[t]{0.31\textwidth}
                \centering
                \includegraphics[width=\linewidth]{04_00_ownership_feature_distribution_03_reported_ownership_hhi.pdf}
        \end{subfigure}
        \par\medskip
        \begin{subfigure}[t]{0.31\textwidth}
                \centering
                \includegraphics[width=\linewidth]{04_00_ownership_feature_distribution_04_top_five_reported_ownership_share.pdf}
        \end{subfigure}
        \hfill
        \begin{subfigure}[t]{0.31\textwidth}
                \centering
                \includegraphics[width=\linewidth]{04_00_ownership_feature_distribution_05_mean_owner_portfolio_weight.pdf}
        \end{subfigure}
        \hfill
        \begin{subfigure}[t]{0.31\textwidth}
                \centering
                \includegraphics[width=\linewidth]{04_00_ownership_feature_distribution_06_maximum_owner_portfolio_weight.pdf}
        \end{subfigure}
        \par\medskip
        \begin{subfigure}[t]{0.31\textwidth}
                \centering
                \includegraphics[width=\linewidth]{04_00_ownership_feature_distribution_07_manager_holder_share.pdf}
        \end{subfigure}
        \hfill
        \begin{subfigure}[t]{0.31\textwidth}
                \centering
                \includegraphics[width=\linewidth]{04_00_ownership_feature_distribution_08_institutional_holder_count_change_percent.pdf}
        \end{subfigure}
        \hfill
        \begin{subfigure}[t]{0.31\textwidth}
                \centering
                \includegraphics[width=\linewidth]{04_00_ownership_feature_distribution_09_total_reported_value_change_percent.pdf}
        \end{subfigure}
        \par\medskip
        \makebox[\textwidth][c]{%
                \begin{subfigure}[t]{0.31\textwidth}
                        \centering
                        \includegraphics[width=\linewidth]{04_00_ownership_feature_distribution_10_reported_ownership_hhi_change.pdf}
                \end{subfigure}%
                \hspace{0.02\textwidth}%
                \begin{subfigure}[t]{0.31\textwidth}
                        \centering
                        \includegraphics[width=\linewidth]{04_00_ownership_feature_distribution_11_top_five_reported_ownership_share_change.pdf}
                \end{subfigure}%
        }
        \caption[Training-sample distributions of ownership predictors]{Training-sample distributions of the eleven conventional ownership predictors. Values are clipped at their 1st and 99th percentiles for display only.}
        \label{fig:ownership-feature-distributions}
\end{figure}

\noindent
Figure~\ref{fig:ownership-temporal-stability} reports the quarterly medians of the conventional ownership predictors over the training period. The observed changes through time support estimating preprocessing parameters from the fitting sample and assessing predictive performance chronologically.

\begin{figure}[H]
        \centering
        \begin{subfigure}[t]{0.31\textwidth}
                \centering
                \includegraphics[width=\linewidth]{04_00_ownership_temporal_stability_01_institutional_holder_count.pdf}
        \end{subfigure}
        \hfill
        \begin{subfigure}[t]{0.31\textwidth}
                \centering
                \includegraphics[width=\linewidth]{04_00_ownership_temporal_stability_02_total_reported_value_usd.pdf}
        \end{subfigure}
        \hfill
        \begin{subfigure}[t]{0.31\textwidth}
                \centering
                \includegraphics[width=\linewidth]{04_00_ownership_temporal_stability_03_reported_ownership_hhi.pdf}
        \end{subfigure}
        \par\medskip
        \begin{subfigure}[t]{0.31\textwidth}
                \centering
                \includegraphics[width=\linewidth]{04_00_ownership_temporal_stability_04_top_five_reported_ownership_share.pdf}
        \end{subfigure}
        \hfill
        \begin{subfigure}[t]{0.31\textwidth}
                \centering
                \includegraphics[width=\linewidth]{04_00_ownership_temporal_stability_05_mean_owner_portfolio_weight.pdf}
        \end{subfigure}
        \hfill
        \begin{subfigure}[t]{0.31\textwidth}
                \centering
                \includegraphics[width=\linewidth]{04_00_ownership_temporal_stability_06_maximum_owner_portfolio_weight.pdf}
        \end{subfigure}
        \par\medskip
        \begin{subfigure}[t]{0.31\textwidth}
                \centering
                \includegraphics[width=\linewidth]{04_00_ownership_temporal_stability_07_manager_holder_share.pdf}
        \end{subfigure}
        \hfill
        \begin{subfigure}[t]{0.31\textwidth}
                \centering
                \includegraphics[width=\linewidth]{04_00_ownership_temporal_stability_08_institutional_holder_count_change_percent.pdf}
        \end{subfigure}
        \hfill
        \begin{subfigure}[t]{0.31\textwidth}
                \centering
                \includegraphics[width=\linewidth]{04_00_ownership_temporal_stability_09_total_reported_value_change_percent.pdf}
        \end{subfigure}
        \par\medskip
        \makebox[\textwidth][c]{%
                \begin{subfigure}[t]{0.31\textwidth}
                        \centering
                        \includegraphics[width=\linewidth]{04_00_ownership_temporal_stability_10_reported_ownership_hhi_change.pdf}
                \end{subfigure}%
                \hspace{0.02\textwidth}%
                \begin{subfigure}[t]{0.31\textwidth}
                        \centering
                        \includegraphics[width=\linewidth]{04_00_ownership_temporal_stability_11_top_five_reported_ownership_share_change.pdf}
                \end{subfigure}%
        }
        \caption[Temporal stability of ownership predictors]{Quarterly medians of the conventional ownership predictors across the training sample.}
        \label{fig:ownership-temporal-stability}
\end{figure}

\section{Ownership Network and Features}
\label{app:ownership_network}

\subsection{Bipartite Manager--Stock Network}
\label{app:bipartite_graph}

Figure~\ref{fig:network-growth} summarizes the growth of the quarterly ownership network in its two node sets and reported ownership edges.

\begin{figure}[H]
        \centering
        \begin{subfigure}[t]{0.48\textwidth}
                \centering
                \includegraphics[width=\linewidth]{02_01_network_nodes_matplotlib.pdf}
                \caption{Manager and security nodes.}
        \end{subfigure}
        \hfill
        \begin{subfigure}[t]{0.48\textwidth}
                \centering
                \includegraphics[width=\linewidth]{02_01_network_edges_matplotlib.pdf}
                \caption{Ownership edges.}
        \end{subfigure}
        \caption{Evolution of the quarterly institutional ownership network.}
        \label{fig:network-growth}
\end{figure}

\noindent
The quarterly bipartite graphs are large but sparse. Their density ranges from 3.16\% to 4.86\% of all possible manager--security pairs, while every quarterly graph forms a single connected component. Thus, any two securities are connected through some chain of shared owners even though most possible ownership edges do not exist. A full quarterly graph is not directly legible because a typical snapshot contains several thousand managers, a comparable number of securities, and hundreds of thousands of ownership edges; the edge count exceeds one million in later quarters.

Figure~\ref{fig:network-degree-distributions} reports the pooled distributions of securities held per manager and institutional holders per stock across the quarterly ownership networks. Both are strongly right-skewed: most nodes have comparatively low degree, while a smaller group of managers holds many securities and a smaller group of stocks is held by many institutions.

\begin{figure}[H]
        \centering
        \begin{subfigure}[t]{0.48\textwidth}
                \centering
                \includegraphics[width=\linewidth]{08_01_manager_degree_distribution.pdf}
                \caption{Securities held per manager.}
                \label{fig:manager-degree-distribution}
        \end{subfigure}
        \hfill
        \begin{subfigure}[t]{0.48\textwidth}
                \centering
                \includegraphics[width=\linewidth]{08_01_stock_degree_distribution.pdf}
                \caption{Institutional holders per stock.}
                \label{fig:stock-degree-distribution}
        \end{subfigure}
        \caption[Node-degree distributions in the ownership network]{Pooled distributions of manager and stock node degrees across the quarterly ownership networks. Counts are shown in their original units on logarithmic horizontal axes.}
        \label{fig:network-degree-distributions}
\end{figure}

\subsection{Manager Similarity and Communities}
\label{app:manager_similarity}

Managers are compared using their quarterly vectors of portfolio weights. For managers \(m\) and \(n\), portfolio similarity is measured by cosine similarity,

\begin{equation}
        \operatorname{sim}_{m,n,q}
        =
        \frac{\mathbf{p}_{m,q}^{\prime}\mathbf{p}_{n,q}}
        {\lVert\mathbf{p}_{m,q}\rVert\lVert\mathbf{p}_{n,q}\rVert},
        \label{eq:manager_cosine_similarity}
\end{equation}

\noindent
where \(\mathbf{p}_{m,q}\) contains manager \(m\)'s portfolio weights in quarter \(q\). Each manager is connected to its 25 nearest positive-similarity neighbours. The resulting relations are converted into an undirected weighted graph, and Louvain community detection is applied separately in each quarter \parencite{blondel2008fast}. Managers in the same community therefore hold relatively similar portfolios within that quarter; community labels are not treated as persistent identities through time.

Figure~\ref{fig:manager-community-evolution} summarizes the resulting community structure. The number of detected communities generally increases, while the share of managers in the largest community declines. This pattern indicates that the expanding manager population becomes divided among a broader set of contemporaneous portfolio-similarity groups.

\begin{figure}[H]
        \centering
        \begin{subfigure}[t]{0.48\textwidth}
                \centering
                \includegraphics[width=\linewidth]{02_02_community_count_matplotlib.pdf}
                \caption{Number of communities.}
        \end{subfigure}
        \hfill
        \begin{subfigure}[t]{0.48\textwidth}
                \centering
                \includegraphics[width=\linewidth]{02_02_largest_community_share_matplotlib.pdf}
                \caption{Largest-community share.}
        \end{subfigure}
        \caption[Quarterly manager-community structure]{Number of manager communities and share of managers in the largest community by reporting quarter.}
        \label{fig:manager-community-evolution}
\end{figure}

\noindent
Across the pooled quarterly sample, the 25-nearest-neighbor manager-similarity graphs contain 5,627,437 edges. Applying Louvain community detection separately in each quarter produces between 7 and 18 communities, with a median of 11.5. Figure~\ref{fig:community-meta-graph} illustrates this structure for the final test quarter. Each detected community is collapsed into a single node, whose size represents its manager count, and the strongest inter-community similarity links are retained for display.

\begin{figure}[H]
        \centering
        \includegraphics[width=0.75\textwidth]{08_01_community_meta_graph.pdf}
        \caption[Manager-community meta-graph]{Manager communities in the illustrative quarter, collapsed to one node per community, with node size representing manager count and the strongest inter-community similarity links displayed.}
        \label{fig:community-meta-graph}
\end{figure}

\subsection{Network Feature Definitions and Diagnostics}
\label{app:network_features}

For each quarter, manager portfolios are represented by vectors of reported portfolio weights. Cosine similarity is calculated between these vectors, and each manager is linked to its 25 nearest positive-similarity neighbours. The directed neighbour relations are converted into an undirected union graph, retaining the larger similarity when a relation is observed in both directions. Louvain communities are estimated on this weighted graph with resolution 1 and random seed 13. Let $a_{m,i,q}=v_{m,i,q}/V_{i,q}$ denote manager $m$'s share of stock $i$'s reported institutional value, $d_{m,q}$ its degree in the manager-similarity graph, $M_q$ the number of managers, and $\bar{s}_{m,q}$ its mean similarity to graph neighbours. Table~\ref{tab:network-feature-definitions} defines the eight resulting stock-level features.

\begingroup
\scriptsize
\setlength{\tabcolsep}{4pt}
\renewcommand{\arraystretch}{1.15}
\begin{longtable}{>{\raggedright\arraybackslash}p{0.24\textwidth}>{\raggedright\arraybackslash}p{0.50\textwidth}>{\raggedright\arraybackslash}p{0.16\textwidth}}
        \caption{Exact definitions of the ownership-network predictors}
        \label{tab:network-feature-definitions}                                                                                                                                                                                             \\
        \toprule
        Feature & Exact construction                                                                                                                                                                                    & Construction rule \\
        \midrule
        \endfirsthead
        \toprule
        Feature & Exact construction                                                                                                                                                                                    & Construction rule \\
        \midrule
        \endhead
        \bottomrule
        \endfoot
        \nolinkurl{owner_similarity_score}
                & Mean cosine similarity across all unordered pairs of the stock's owners. Owner pairs not connected in the 25-nearest-neighbour graph contribute zero; stocks with fewer than two owners receive zero.
                & Calculated in levels.                                                                                                                                                                                                     \\
        \nolinkurl{value_weighted_owner_degree_centrality}
                & $\sum_m a_{m,i,q}\,d_{m,q}/(M_q-1)$.
                & Weighted by reported position value.                                                                                                                                                                                      \\
        \nolinkurl{value_weighted_owner_neighbor_similarity}
                & $\sum_m a_{m,i,q}\,\bar{s}_{m,q}$.
                & Weighted by reported position value.                                                                                                                                                                                      \\
        \nolinkurl{owner_community_count}
                & Number of distinct contemporaneous Louvain communities represented among the stock's owners.
                & Calculated in levels.                                                                                                                                                                                                     \\
        \nolinkurl{owner_community_hhi}
                & $\sum_c A_{c,i,q}^2$, where $A_{c,i,q}=\sum_{m\in c}a_{m,i,q}$ is community $c$'s share of the stock's reported institutional value.
                & Calculated in levels.                                                                                                                                                                                                     \\
        \nolinkurl{largest_owner_community_share}
                & $\max_c A_{c,i,q}$.
                & Calculated in levels.                                                                                                                                                                                                     \\
        \nolinkurl{owner_similarity_score_change}
                & Difference between the current and preceding quarter's owner-similarity score.
                & Defined only for consecutive observed quarters.                                                                                                                                                                           \\
        \nolinkurl{owner_community_hhi_change}
                & Difference between the current and preceding quarter's owner-community HHI.
                & Defined only for consecutive observed quarters.                                                                                                                                                                           \\
\end{longtable}
\endgroup

\noindent
The level features require at least one retained owner, except that owner similarity is set to zero for stocks with fewer than two owners. Change features require observations in consecutive reporting quarters.

Raw network-feature levels can shift mechanically with the size and composition of the graph. Figure~\ref{fig:network-temporal-stability} tracks the quarterly median of each raw feature alongside the number of reporting managers. Manager count correlates $-0.922$ with median owner similarity, $-0.855$ with degree centrality, $0.803$ with community count, and $0.865$ with neighbor similarity. Because the manager population grows substantially over the sample, these relationships indicate that raw feature levels partly drift for mechanical rather than economic reasons and motivate the within-quarter percentile-rank transformation used in the predictive analysis.

\begin{figure}[H]
        \centering
        \begin{subfigure}[t]{0.48\textwidth}
                \centering
                \includegraphics[width=\linewidth]{02_04_temporal_stability_01_owner_similarity_score.pdf}
        \end{subfigure}
        \hfill
        \begin{subfigure}[t]{0.48\textwidth}
                \centering
                \includegraphics[width=\linewidth]{02_04_temporal_stability_02_value_weighted_owner_degree_centrality.pdf}
        \end{subfigure}
        \par\medskip
        \begin{subfigure}[t]{0.48\textwidth}
                \centering
                \includegraphics[width=\linewidth]{02_04_temporal_stability_03_value_weighted_owner_neighbor_similarity.pdf}
        \end{subfigure}
        \hfill
        \begin{subfigure}[t]{0.48\textwidth}
                \centering
                \includegraphics[width=\linewidth]{02_04_temporal_stability_04_owner_community_count.pdf}
        \end{subfigure}
        \par\medskip
        \begin{subfigure}[t]{0.48\textwidth}
                \centering
                \includegraphics[width=\linewidth]{02_04_temporal_stability_05_owner_community_hhi.pdf}
        \end{subfigure}
        \hfill
        \begin{subfigure}[t]{0.48\textwidth}
                \centering
                \includegraphics[width=\linewidth]{02_04_temporal_stability_06_largest_owner_community_share.pdf}
        \end{subfigure}
        \par\medskip
        \begin{subfigure}[t]{0.48\textwidth}
                \centering
                \includegraphics[width=\linewidth]{02_04_temporal_stability_07_owner_similarity_score_change.pdf}
        \end{subfigure}
        \hfill
        \begin{subfigure}[t]{0.48\textwidth}
                \centering
                \includegraphics[width=\linewidth]{02_04_temporal_stability_08_owner_community_hhi_change.pdf}
        \end{subfigure}
        \caption[Temporal stability of raw network features]{Quarterly median of each raw network feature across the full sample period, before the within-quarter percentile-rank transformation used in modeling.}
        \label{fig:network-temporal-stability}
\end{figure}

\noindent
Figure~\ref{fig:app-network-distributions} reports the underlying training-sample distributions of the eight network features introduced in Section~\ref{sec:network_features}. Correlations involving these features are evaluated jointly with the market and ownership features in Section~\ref{sec:feature_engineering}.

\begin{figure}[H]
        \centering
        \begin{subfigure}[t]{0.48\textwidth}
                \centering
                \includegraphics[width=\linewidth]{02_04_feature_distribution_01_owner_similarity_score.pdf}
        \end{subfigure}
        \hfill
        \begin{subfigure}[t]{0.48\textwidth}
                \centering
                \includegraphics[width=\linewidth]{02_04_feature_distribution_02_value_weighted_owner_degree_centrality.pdf}
        \end{subfigure}
        \par\medskip
        \begin{subfigure}[t]{0.48\textwidth}
                \centering
                \includegraphics[width=\linewidth]{02_04_feature_distribution_03_value_weighted_owner_neighbor_similarity.pdf}
        \end{subfigure}
        \hfill
        \begin{subfigure}[t]{0.48\textwidth}
                \centering
                \includegraphics[width=\linewidth]{02_04_feature_distribution_04_owner_community_count.pdf}
        \end{subfigure}
        \par\medskip
        \begin{subfigure}[t]{0.48\textwidth}
                \centering
                \includegraphics[width=\linewidth]{02_04_feature_distribution_05_owner_community_hhi.pdf}
        \end{subfigure}
        \hfill
        \begin{subfigure}[t]{0.48\textwidth}
                \centering
                \includegraphics[width=\linewidth]{02_04_feature_distribution_06_largest_owner_community_share.pdf}
        \end{subfigure}
        \par\medskip
        \begin{subfigure}[t]{0.48\textwidth}
                \centering
                \includegraphics[width=\linewidth]{02_04_feature_distribution_07_owner_similarity_score_change.pdf}
        \end{subfigure}
        \hfill
        \begin{subfigure}[t]{0.48\textwidth}
                \centering
                \includegraphics[width=\linewidth]{02_04_feature_distribution_08_owner_community_hhi_change.pdf}
        \end{subfigure}
        \caption[Training-sample network-feature distributions]{Distribution of each of the eight raw network features in the training sample, capped at the 1st and 99th percentiles for display only.}
        \label{fig:app-network-distributions}
\end{figure}

\section{Feature Processing and Selection}
\label{app:feature_processing}

\subsection{Feature Missingness}

Table~\ref{tab:feature-missingness} reports the share of missing observations for every candidate predictor before imputation and redundancy-based feature selection. All market features and all ownership and network level features are complete in the three modeling samples. Missingness is confined to the four ownership-change measures and two network-change measures, which require the stock to be observed in consecutive reporting quarters. Their missing share is higher in the training sample because the beginning of the panel contains no preceding-quarter history, while only a small share of later validation and test observations lacks a consecutive-quarter predecessor.

\begingroup
\footnotesize
\setlength{\tabcolsep}{4pt}
\renewcommand{\arraystretch}{1.08}
\begin{longtable}{p{0.14\textwidth}p{0.40\textwidth}rrr}
    \caption{Missingness across candidate predictors and modeling samples}
    \label{tab:feature-missingness} \\
    \toprule
    Feature group & Feature & \multicolumn{3}{c}{Missing share} \\
    \cmidrule(lr){3-5}
                  &         & Training & Validation & Test \\
    \midrule
    \endfirsthead
    \multicolumn{5}{l}{\footnotesize\itshape Table~\thetable\ continued} \\
    \toprule
    Feature group & Feature & \multicolumn{3}{c}{Missing share} \\
    \cmidrule(lr){3-5}
                  &         & Training & Validation & Test \\
    \midrule
    \endhead
    \midrule
    \multicolumn{5}{r}{\footnotesize\itshape Continued on next page} \\
    \endfoot
    \bottomrule
    \endlastfoot
        Market & Market capitalization & 0.00\% & 0.00\% & 0.00\% \\
        Market & Average daily dollar volume & 0.00\% & 0.00\% & 0.00\% \\
        Market & 21-day return & 0.00\% & 0.00\% & 0.00\% \\
        Market & Skip-month momentum & 0.00\% & 0.00\% & 0.00\% \\
        Market & Volatility & 0.00\% & 0.00\% & 0.00\% \\
        Market & Downside volatility & 0.00\% & 0.00\% & 0.00\% \\
        Market & Market beta & 0.00\% & 0.00\% & 0.00\% \\
        Market & Past maximum drawdown & 0.00\% & 0.00\% & 0.00\% \\
        Market & Negative-return share & 0.00\% & 0.00\% & 0.00\% \\
        Ownership & Institutional holder count & 0.00\% & 0.00\% & 0.00\% \\
        Ownership & Total reported institutional value & 0.00\% & 0.00\% & 0.00\% \\
        Ownership & Ownership HHI & 0.00\% & 0.00\% & 0.00\% \\
        Ownership & Top-five ownership share & 0.00\% & 0.00\% & 0.00\% \\
        Ownership & Mean owner portfolio weight & 0.00\% & 0.00\% & 0.00\% \\
        Ownership & Maximum owner portfolio weight & 0.00\% & 0.00\% & 0.00\% \\
        Ownership & Manager-holder share & 0.00\% & 0.00\% & 0.00\% \\
        Ownership & Change in institutional holder count & 3.02\% & 0.16\% & 0.15\% \\
        Ownership & Change in total reported institutional value & 3.02\% & 0.16\% & 0.15\% \\
        Ownership & Change in ownership HHI & 3.02\% & 0.16\% & 0.15\% \\
        Ownership & Change in top-five ownership share & 3.02\% & 0.16\% & 0.15\% \\
        Network & Owner similarity & 0.00\% & 0.00\% & 0.00\% \\
        Network & Owner degree centrality & 0.00\% & 0.00\% & 0.00\% \\
        Network & Owner-neighbour similarity & 0.00\% & 0.00\% & 0.00\% \\
        Network & Owner community count & 0.00\% & 0.00\% & 0.00\% \\
        Network & Owner-community HHI & 0.00\% & 0.00\% & 0.00\% \\
        Network & Largest owner-community share & 0.00\% & 0.00\% & 0.00\% \\
        Network & Change in owner similarity & 3.02\% & 0.16\% & 0.15\% \\
        Network & Change in owner-community HHI & 3.02\% & 0.16\% & 0.15\% \\
\end{longtable}
\endgroup


\noindent
Market capitalization, average daily dollar volume, institutional holder count, and total reported institutional value are examined under $\log(1+x)$ transformations because their raw distributions are strongly right-skewed. Figure~\ref{fig:log-transformations} shows the two market scale variables, while Figure~\ref{fig:ownership-log-transformations} presents the corresponding ownership variables. After redundancy-based selection, only market capitalization remains in the fitted tabular and stock-node feature sets. Raw distributions are capped at their 99th percentiles for display only.

\begin{figure}[H]
        \centering
        \begin{subfigure}[t]{0.48\textwidth}
                \centering
                \includegraphics[width=\linewidth]{04_00_market_log_transformation_01_market_cap_usd_raw.pdf}
        \end{subfigure}
        \hfill
        \begin{subfigure}[t]{0.48\textwidth}
                \centering
                \includegraphics[width=\linewidth]{04_00_market_log_transformation_02_market_cap_usd_log1p.pdf}
        \end{subfigure}
        \par\medskip
        \begin{subfigure}[t]{0.48\textwidth}
                \centering
                \includegraphics[width=\linewidth]{04_00_market_log_transformation_03_average_daily_dollar_volume_63d_raw.pdf}
        \end{subfigure}
        \hfill
        \begin{subfigure}[t]{0.48\textwidth}
                \centering
                \includegraphics[width=\linewidth]{04_00_market_log_transformation_04_average_daily_dollar_volume_63d_log1p.pdf}
        \end{subfigure}
        \caption[Logarithmic transformations of market scale variables]{Raw and $\log(1+x)$-transformed distributions of market capitalization and average daily dollar volume. Raw panels are capped at their 99th percentiles for display only.}
        \label{fig:log-transformations}
\end{figure}

\begin{figure}[H]
        \centering
        \begin{subfigure}[t]{0.48\textwidth}
                \centering
                \includegraphics[width=\linewidth]{04_00_ownership_log_transformation_01_institutional_holder_count_raw.pdf}
        \end{subfigure}
        \hfill
        \begin{subfigure}[t]{0.48\textwidth}
                \centering
                \includegraphics[width=\linewidth]{04_00_ownership_log_transformation_02_institutional_holder_count_log1p.pdf}
        \end{subfigure}
        \par\medskip
        \begin{subfigure}[t]{0.48\textwidth}
                \centering
                \includegraphics[width=\linewidth]{04_00_ownership_log_transformation_03_total_reported_value_usd_raw.pdf}
        \end{subfigure}
        \hfill
        \begin{subfigure}[t]{0.48\textwidth}
                \centering
                \includegraphics[width=\linewidth]{04_00_ownership_log_transformation_04_total_reported_value_usd_log1p.pdf}
        \end{subfigure}
        \caption[Logarithmic transformations of ownership scale variables]{Raw and $\log(1+x)$-transformed distributions of institutional holder count and total reported institutional value. Raw panels are capped at their 99th percentiles for display only.}
        \label{fig:ownership-log-transformations}
\end{figure}

\noindent
Each retained tabular predictor and graph-node feature is clipped at the 1st and 99th percentiles estimated from the relevant fitting sample, and remaining missing values are replaced by fitting-sample medians. Standardization is applied where required by the model class. The fitted preprocessing parameters are applied unchanged outside the fitting sample, preventing validation or test information from influencing the transformations, clipping bounds, imputation values, or scaling moments.

Before modeling, each network feature is converted to an average-rank percentile within its full quarterly stock universe and scaled from 0 to 1. A feature receives 0.5 if its quarter contains only one valid observation. Ranking is performed before modeling-eligibility restrictions are applied, while missing change features are subsequently handled by fitting-sample median imputation.

\subsection{Strong Feature Correlations}
\label{app:feature_correlations}

Table~\ref{tab:app-strong-feature-correlations} lists every unique feature pair whose absolute mean quarterly Spearman correlation is at least 0.70 in the training sample. Correlations are calculated separately within each quarter and then averaged across quarters, giving every training quarter equal weight. The threshold is used as a compact redundancy diagnostic rather than as an automatic variable-selection rule.

\begingroup
\footnotesize
\setlength{\tabcolsep}{5pt}
\begin{longtable}{>{\raggedright\arraybackslash}p{0.57\textwidth}>{\raggedright\arraybackslash}p{0.25\textwidth}r}
        \caption{Feature pairs with absolute mean quarterly Spearman correlation of at least 0.70}
        \label{tab:app-strong-feature-correlations}                                    \\
        \toprule
        Feature pair                                   & Groups               & $\rho$ \\
        \midrule
        \endfirsthead
        \multicolumn{3}{l}{\footnotesize\itshape Table~\thetable\ continued}           \\
        \toprule
        Feature pair                                   & Groups               & $\rho$ \\
        \midrule
        \endhead
        \midrule
        \multicolumn{3}{r}{\footnotesize\itshape Continued on next page}               \\
        \endfoot
        \bottomrule
        \endlastfoot
        Holder count -- manager-holder share           & Ownership--ownership & 1.000  \\
        Ownership HHI -- top-five ownership share      & Ownership--ownership & 0.963  \\
        Owner-community HHI -- largest community share & Network--network     & 0.958  \\
        Market capitalization -- holder count          & Market--ownership    & 0.954  \\
        Market capitalization -- manager-holder share  & Market--ownership    & 0.954  \\
        Market capitalization -- reported value        & Market--ownership    & 0.943  \\
        Reported value -- manager-holder share         & Ownership--ownership & 0.938  \\
        Holder count -- reported value                 & Ownership--ownership & 0.938  \\
        Dollar volume -- holder count                  & Market--ownership    & 0.937  \\
        Dollar volume -- manager-holder share          & Market--ownership    & 0.937  \\
        Volatility -- downside volatility              & Market--market       & 0.936  \\
        Market capitalization -- dollar volume         & Market--market       & 0.926  \\
        Mean owner weight -- maximum owner weight      & Ownership--ownership & 0.902  \\
        Dollar volume -- reported value                & Market--ownership    & 0.900  \\
        Holder count -- owner-community count          & Ownership--network   & 0.891  \\
        Manager-holder share -- owner-community count  & Ownership--network   & 0.891  \\
        Market capitalization -- owner-community count & Market--network      & 0.844  \\
        Dollar volume -- owner-community count         & Market--network      & 0.834  \\
        Reported value -- owner-community count        & Ownership--network   & 0.831  \\
        Ownership-HHI change -- top-five-share change  & Ownership--ownership & 0.806  \\
        Volatility -- past maximum drawdown            & Market--market       & 0.787  \\
        Downside volatility -- past maximum drawdown   & Market--market       & 0.785  \\
\end{longtable}
\endgroup

\subsection{Redundancy-Based Feature Selection}
\label{app:feature_selection_decisions}

Table~\ref{tab:correlation-feature-selection} summarizes the feature-selection decisions produced by applying the correlation threshold and the interpretability hierarchy described in Section~\ref{sec:feature_engineering}.

\begin{table}[H]
        \centering
        \footnotesize
        \caption{Redundancy-based feature-selection decisions}
        \label{tab:correlation-feature-selection}
        \begin{tabular}{>{\raggedright\arraybackslash}p{0.22\textwidth}>{\raggedright\arraybackslash}p{0.31\textwidth}>{\raggedright\arraybackslash}p{0.36\textwidth}}
                \toprule
                Retained & Excluded                                                                                                                    & Reason \\
                \midrule
                Market capitalization
                         & Dollar volume; holder count; reported value; manager-holder share; owner-community count
                         & Retains the standard measure of firm scale and removes closely related measures of trading and institutional participation.          \\
                Downside volatility
                         & Total volatility
                         & Measures adverse return variation and is more closely aligned with the prediction target.                                            \\
                Ownership HHI
                         & Top-five ownership share
                         & Uses the complete ownership distribution rather than only its largest positions.                                                     \\
                Mean owner portfolio weight
                         & Maximum owner portfolio weight
                         & Represents the stock's importance across all owners rather than one extreme owner.                                                   \\
                Change in ownership HHI
                         & Change in top-five ownership share
                         & Preserves consistency with the retained concentration level.                                                                         \\
                Owner-community HHI
                         & Largest owner-community share
                         & Uses the complete distribution of ownership across communities.                                                                      \\
                \bottomrule
        \end{tabular}
\end{table}

\section{Chronological Sample Allocation}
\label{app:chronological_samples}

The chronological allocation excludes otherwise eligible observations from 2021~Q4 and 2023~Q4. These quarters separate the training sample from validation and validation from test, respectively. At each boundary, the latest target window in the preceding sample closes before the earliest information cutoff in the following sample, ensuring that outcomes used in one sample are fully resolved before the next sample begins.

Figure~\ref{fig:sample-target-distribution} compares the realized distribution of the primary target across the training, validation, and test samples. It documents a rightward shift after the training period, which the fixed chronological evaluation must confront out of sample.

\begin{figure}[H]
        \centering
        \includegraphics[width=0.75\textwidth]{01_08_target_distribution_boxplot.pdf}
        \caption[Target distribution across the modeling samples]{Distribution of the forward 63-day maximum drawdown target within each chronological sample. Outlier markers are hidden for readability only; all observations are included in the boxplot statistics.}
        \label{fig:sample-target-distribution}
\end{figure}

\section{Graph Neural Network Architectures}
\label{app:gnn_architectures}

All three graph neural networks operate on the quarterly bipartite manager--stock graphs defined in Section~\ref{sec:bipartite_graph}. Each architecture first projects the stock and manager input variables into a common hidden dimension and then passes messages from stocks to managers and from managers back to stocks. The resulting stock representation is mapped to a prediction of forward downside risk. The architectures differ in how neighboring nodes are aggregated and in whether information is carried across quarters.

\subsection{Weighted GraphSAGE}

The GraphSAGE specification uses observed holdings weights to aggregate neighboring representations. Let \(\mathbf{h}^{(0)}_{i,q}\) and \(\mathbf{h}^{(0)}_{m,q}\) denote the initial hidden representations of stock \(i\) and manager \(m\) in quarter \(q\). In the stock-to-manager step, stock messages are weighted by their portfolio weights \(p_{mi,q}\):

\begin{equation}
    \mathbf{h}^{(1)}_{m,q}
    =
    \operatorname{ReLU}
    \left(
        W^{\mathrm{self}}_M\mathbf{h}^{(0)}_{m,q}
        +
        W^{\mathrm{nbr}}_M
        \sum_{i\in\mathcal{N}_q(m)}
        p_{mi,q}\mathbf{h}^{(0)}_{i,q}
    \right).
    \label{eq:weighted_graphsage_manager}
\end{equation}

\noindent
The manager representations are then passed back to stocks using each manager's share \(o_{mi,q}\) of the stock's reported institutional ownership:

\begin{equation}
    \mathbf{g}_{i,q}
    =
    \operatorname{Dropout}
    \left[
    \operatorname{ReLU}
    \left(
        W^{\mathrm{self}}_S\mathbf{h}^{(0)}_{i,q}
        +
        W^{\mathrm{nbr}}_S
        \sum_{m\in\mathcal{N}_q(i)}
        o_{mi,q}\mathbf{h}^{(1)}_{m,q}
    \right)
    \right].
    \label{eq:weighted_graphsage_stock}
\end{equation}

\noindent
The static GraphSAGE prediction is \(\hat{y}_{i,q}=\mathbf{w}_o^{\prime}\mathbf{g}_{i,q}+b_o\). The two sets of observed edge weights determine how strongly each neighbor contributes; the model learns the transformations applied to a node's own representation and to the weighted neighbor aggregate. Each quarterly graph is processed independently.

\subsection{Graph Attention Network}

GAT uses the same two-direction message-passing sequence but learns the relative importance of adjacent nodes. For an edge from source node \(u\) to destination node \(v\), an unnormalized attention score can be represented as

\begin{equation}
    e_{uv,q}
    =
    \operatorname{LeakyReLU}
    \left[
        \mathbf{a}^{\prime}
        \left(
            \Theta_s\mathbf{h}_{u,q}
            \,\Vert\,
            \Theta_d\mathbf{h}_{v,q}
            \,\Vert\,
            \Theta_e\mathbf{e}_{uv,q}
        \right)
    \right],
    \label{eq:gat_attention_score}
\end{equation}

\noindent
where \(\mathbf{e}_{uv,q}\) contains the standardized logarithms of portfolio weight and reported-ownership share. Scores are normalized across the neighbors of the destination node,

\begin{equation}
    \alpha_{uv,q}
    =
    \frac{\exp(e_{uv,q})}
    {\sum_{k\in\mathcal{N}_q(v)}\exp(e_{kv,q})},
    \qquad
    \mathbf{m}_{v,q}
    =
    \sum_{u\in\mathcal{N}_q(v)}
    \alpha_{uv,q}\Theta_s\mathbf{h}_{u,q}.
    \label{eq:gat_attention_weights}
\end{equation}

\noindent
The implementation uses one attention head in each bipartite direction and combines the resulting message with the destination node through a residual update before applying ReLU. Unlike weighted GraphSAGE, GAT is not constrained to treat the reported ownership weights as the final importance assigned to neighboring managers or stocks. It can learn different attention weights from the joint node and edge information, but this additional flexibility also introduces more parameters. Each quarter is again processed independently.

\subsection{Temporal GraphSAGE}

Temporal GraphSAGE first obtains the contemporaneous stock embedding \(\mathbf{g}_{i,q}\) from the weighted GraphSAGE encoder. A gated recurrent unit then combines it with the stock's state from the preceding quarter:

\begin{equation}
    \mathbf{s}_{i,q}
    =
    \operatorname{GRU}
    \left(
        \mathbf{g}_{i,q},
        \mathbf{s}_{i,q-1}
    \right),
    \qquad
    \hat{y}_{i,q}
    =
    \mathbf{w}_o^{\prime}\mathbf{s}_{i,q}+b_o.
    \label{eq:temporal_graphsage}
\end{equation}

\noindent
This architecture therefore combines the current ownership network with information retained from previous quarterly snapshots. The preceding state is reset to zero when a stock was absent from the immediately preceding quarter. Purged snapshots update the recurrent state but do not contribute to the training loss.

In summary, weighted GraphSAGE aggregates current neighbors using observed ownership weights, GAT learns the relative importance of current neighbors, and temporal GraphSAGE adds memory of how each stock's graph representation evolves through time.

\section{Model Settings and Selection Rules}
\label{app:model_specifications}

Table~\ref{tab:model-specifications} reports the fixed model settings, controlled validation-stage refinements, and staged selection rules used in the predictive analysis. The common project seed is 42. Ridge is deterministic under the stated specification; the seed is applied to XGBoost and to the Python, NumPy, and PyTorch random-number generators used for graph-model estimation. PyTorch deterministic algorithms are enabled for the GNNs.

\begingroup
\scriptsize
\setlength{\tabcolsep}{4pt}
\renewcommand{\arraystretch}{1.15}
\begin{longtable}{>{\raggedright\arraybackslash}p{0.19\textwidth}>{\raggedright\arraybackslash}p{0.24\textwidth}>{\raggedright\arraybackslash}p{0.47\textwidth}}
        \caption{Predictive-model specifications and selection rules}
        \label{tab:model-specifications}                                                                                                                                                                                                                                                                                                                                                                                                                                                                                                          \\
        \toprule
        Model & Component                                                                                                                                                                                                                                                                                                                                                                                                                                                                                                         & Specification \\
        \midrule
        \endfirsthead
        \toprule
        Model & Component                                                                                                                                                                                                                                                                                                                                                                                                                                                                                                         & Specification \\
        \midrule
        \endhead
        \bottomrule
        \endfoot
        No-information benchmark
              & Forecast
              & Constant equal to the fitting-sample mean target: the training mean for validation and the combined training--validation mean for the final test comparison.                                                                                                                                                                                                                                                                                                                                                                      \\
        All fitted models
              & Selection rule
              & Validation MAE is the primary quantitative criterion. When MAE differences are negligible, RMSE, $R^2$, mean quarterly Spearman correlation, convergence, and model complexity are used to prefer the more stable and parsimonious specification.                                                                                                                                                                                                                                                                             \\
        Modeling ladder
              & Staged selection
              & First, the three nested feature sets are compared separately within Ridge and XGBoost under common model-specific settings. One feature set is retained within each class, using parsimony when validation performance is materially similar. Second, the retained XGBoost specification is refined through a controlled comparison while Ridge remains fixed. Third, the three GNN architectures are compared under a common configuration and only the strongest architecture is refined. The retained Ridge, XGBoost, and GNN specifications then enter the final validation comparison.                                                       \\
        Finalists
              & Refitting
              & The feature sets, hyperparameters, and validation-selected number of boosting rounds or training epochs are held fixed while all three finalists are re-estimated on the combined training and validation samples. They are evaluated once on the test sample to assess generalization without reopening model selection.                                                                                                                                                                                                      \\
        Ridge
              & Estimator
              & Linear regression with an $L_2$ penalty; predictors are standardized using fitting-sample moments.                                                                                                                                                                                                                                                                                                                                                                                                                              \\
        Ridge
              & Penalty
              & All feature-set and finalist comparisons use the prespecified value $\alpha=1$; no separate Ridge tuning stage is performed.                                                                                                                                                                                                                                                                                                                                                                                                 \\
        XGBoost
              & Training objective
              & Squared-error regression objective (\texttt{reg:squarederror}); validation MAE is the primary evaluation and selection metric. Trees are constructed using the histogram method.                                                                                                                                                                                                                                                                                                                                                  \\
        XGBoost
              & Initial settings
              & The feature-set comparison permits at most 1,000 boosting rounds and uses maximum depth 3, minimum child weight 1, row subsampling 0.8, learning rate 0.05, and column subsampling 0.8. Validation MAE is monitored with early-stopping patience of 30 rounds.                                                                                                                                                                                                                                                                         \\
        XGBoost
              & Controlled refinement
              & The retained specification is compared under three configurations: $(\eta,\text{depth})\in\{(0.05,3),(0.10,3),(0.10,4)\}$. Minimum child weight remains 1, row and column subsampling remain 0.8, and each configuration permits at most 1,000 rounds with early-stopping patience of 30. The depth-three, $\eta=0.10$ configuration is retained at its validation-selected round because it matches the validation performance of the alternatives with faster convergence and less complexity than the deeper model. \\
        All GNNs
              & Node and edge preprocessing
              & Stock nodes use the seven retained market and five retained ownership features. Manager nodes use portfolio value, security count, largest position weight, and portfolio HHI. Market capitalization, manager portfolio value, and manager security count receive $\log(1+x)$ transformations. Node features are clipped, median-imputed, and standardized using the fitting sample. For GAT, the logarithms of portfolio weight and reported-ownership share are standardized using fitting-sample edge moments.                 \\
        All GNNs
              & Initial comparison
              & GraphSAGE, GAT, and temporal GraphSAGE are first compared using a common hidden dimension of 16, dropout of 0.20, learning rate of 0.001, and weight decay of 0.0001. Each architecture uses one stock-to-manager and one manager-to-stock message-passing stage, ReLU activation, and a scalar linear output layer.                                                                                                                                                                                                                  \\
        All GNNs
              & Estimation
              & $L_1$ training loss and Adam optimizer; at most 50 epochs; validation MAE evaluated after every epoch; early-stopping patience 5 with minimum improvement 0.0001. The checkpoint with the lowest validation MAE is retained.                                                                                                                                                                                                                                                                                                      \\
        Selected GNN
              & Controlled refinement
              & The retained Temporal GraphSAGE architecture is evaluated under the base configuration $(16,0.20,0.001)$ and four variants $(16,0.20,0.010)$, $(32,0.20,0.001)$, $(32,0.30,0.001)$, and $(64,0.30,0.001)$, where each tuple gives hidden dimension, dropout, and learning rate. Weight decay remains 0.0001 and the same early-stopping rule is used. The configuration $(32,0.20,0.001)$ is retained at epoch 28.                                                                                                                         \\
        GraphSAGE
              & Aggregation
              & Weighted neighbour aggregation in both bipartite directions. Stock messages entering a manager are weighted by portfolio weight; manager messages entering a stock are weighted by reported-ownership share. Each update combines separate linear transformations of the node's current representation and its weighted neighbour aggregate.                                                                                                                                                                                      \\
        GAT
              & Attention
              & Two single-head bipartite attention layers, one in each direction, with no self-loops, no concatenation across heads, residual updates, and two edge attributes: the standardized logarithms of portfolio weight and reported-ownership share.                                                                                                                                                                                                                                                                                  \\
        Temporal GraphSAGE
              & Recurrent state
              & The contemporaneous GraphSAGE stock embedding and the stock's preceding-quarter hidden state enter a gated recurrent unit. The prior state is reset to zero when the stock was absent from the immediately preceding quarter. Purged snapshots update states but contribute no training loss.                                                                                                                                                                                                                                     \\
\end{longtable}
\endgroup
